This appendix tabulates the calibration-accuracy metrics and the experiment-marginalised nuisance priors for the fiducial \texttt{Titanium}~+~\texttt{Hybrid} configuration, complementing the visual summaries presented in Sections~\ref{sec:4} and~\ref{sec:5}. Table~\ref{tab:A1} lists, for each tomographic bin \(m\), the within-experiment calibration-accuracy metrics evaluated across the \(501\) realised statistical experiments (Section~\ref{sec:4.2}) and the deterministic correction parameters and one-sigma priors derived from the Dirichlet-sampled Monte Carlo ensemble (Sections~\ref{sec:5.2} and~\ref{sec:5.3}). These values are intended as a compact numerical reference for downstream forecasting studies that adopt our calibration outputs.

\begin{table*}
  \centering
  \footnotesize
  \begin{threeparttable}
    \caption{Summary of the calibration-accuracy metrics and experiment-marginalised nuisance priors for the tomographic redshift distributions \(\phi_m(z)\) under the fiducial \texttt{Titanium}~+~\texttt{Hybrid} configuration. The left block summarises the scaled within-experiment biases across the ensemble of \(501\) realised statistical experiments (Equation~\ref{eq:22}; Section~\ref{sec:4.2}), quantifying within-experiment calibration accuracy relative to the supported \texttt{Truth} benchmark. The right block lists the deterministic correction parameters \(\langle\Delta_m\rangle\), \(\langle\Gamma_m\rangle\) and their one-sigma residual scatters \(\varsigma_{\Delta_m}\), \(\varsigma_{\Gamma_m}\) evaluated across the Dirichlet-sampled Monte Carlo ensemble (Equations~\ref{eq:34}--\ref{eq:35}; Section~\ref{sec:5.3}), which together define the per-bin offsets and nuisance priors adopted in the released bias-corrected products. Unlike the scaled biases in the left block, \(\langle\Delta_m\rangle\) is quoted directly in redshift units without any \((1+\mu)\) normalisation, and \(\langle\Gamma_m\rangle\) is a multiplicative width-rescaling factor that equals unity in the absence of correction.}
    \label{tab:A1}
    \begin{tabular}{llccccccccc}
      \toprule
      & & & \multicolumn{4}{c}{Within-experiment calibration accuracy} & \multicolumn{4}{c}{Experiment-marginalised priors and offsets} \\ 
      \cmidrule(lr){4-7} \cmidrule(lr){8-11}
      Scenario & Sample & Bin & \(\bar{\delta}_{\mu_m}\) & \(\sigma_{\mu_m}\) & \(\bar{\delta}_{\eta_m}\) & \(\sigma_{\eta_m}\) & \(\langle \Delta_m \rangle\) & \(\varsigma_{\Delta_m}\) & \(\langle \Gamma_m \rangle\) & \(\varsigma_{\Gamma_m}\) \\
      \midrule
      LSST-Y1 & \texttt{lens} & \(m=1\) & +0.001 & +0.001 & -0.000 & +0.001 & -0.002 & +0.001 & +1.002 & +0.010 \\
      LSST-Y1 & \texttt{lens} & \(m=2\) & +0.002 & +0.001 & -0.001 & +0.001 & -0.002 & +0.001 & +1.035 & +0.012 \\
      LSST-Y1 & \texttt{lens} & \(m=3\) & +0.000 & +0.001 & -0.001 & +0.001 & -0.001 & +0.001 & +1.029 & +0.013 \\
      LSST-Y1 & \texttt{lens} & \(m=4\) & -0.001 & +0.002 & -0.002 & +0.001 & +0.002 & +0.002 & +1.050 & +0.017 \\
      LSST-Y1 & \texttt{lens} & \(m=5\) & -0.009 & +0.007 & -0.006 & +0.002 & +0.020 & +0.008 & +1.207 & +0.049 \\
      \midrule
      LSST-Y1 & \texttt{source} & \(m=1\) & -0.000 & +0.002 & -0.008 & +0.003 & +0.000 & +0.010 & +1.102 & +0.044 \\
      LSST-Y1 & \texttt{source} & \(m=2\) & +0.002 & +0.001 & +0.002 & +0.001 & -0.003 & +0.015 & +0.989 & +0.074 \\
      LSST-Y1 & \texttt{source} & \(m=3\) & +0.001 & +0.001 & +0.005 & +0.003 & -0.002 & +0.004 & +0.892 & +0.053 \\
      LSST-Y1 & \texttt{source} & \(m=4\) & +0.001 & +0.002 & -0.000 & +0.001 & -0.001 & +0.004 & +1.011 & +0.040 \\
      LSST-Y1 & \texttt{source} & \(m=5\) & -0.016 & +0.006 & -0.017 & +0.004 & +0.034 & +0.010 & +1.179 & +0.032 \\
      \midrule
      LSST-Y10 & \texttt{lens} & \(m=1\)  & +0.002 & +0.001 & -0.001 & +0.000 & -0.002 & +0.001 & +1.029 & +0.009 \\
      LSST-Y10 & \texttt{lens} & \(m=2\)  & +0.002 & +0.001 & -0.001 & +0.000 & -0.003 & +0.001 & +1.055 & +0.010 \\
      LSST-Y10 & \texttt{lens} & \(m=3\)  & +0.003 & +0.001 & -0.001 & +0.000 & -0.004 & +0.001 & +1.046 & +0.009 \\
      LSST-Y10 & \texttt{lens} & \(m=4\)  & +0.002 & +0.001 & -0.000 & +0.000 & -0.003 & +0.001 & +1.025 & +0.009 \\
      LSST-Y10 & \texttt{lens} & \(m=5\)  & +0.001 & +0.001 & -0.001 & +0.000 & -0.002 & +0.001 & +1.058 & +0.009 \\
      LSST-Y10 & \texttt{lens} & \(m=6\)  & +0.001 & +0.000 & -0.000 & +0.000 & -0.001 & +0.001 & +1.002 & +0.007 \\
      LSST-Y10 & \texttt{lens} & \(m=7\)  & +0.001 & +0.000 & -0.001 & +0.000 & -0.001 & +0.001 & +1.060 & +0.009 \\
      LSST-Y10 & \texttt{lens} & \(m=8\)  & +0.001 & +0.001 & -0.000 & +0.000 & -0.001 & +0.001 & +1.030 & +0.010 \\
      LSST-Y10 & \texttt{lens} & \(m=9\)  & -0.001 & +0.002 & -0.001 & +0.001 & +0.001 & +0.001 & +1.088 & +0.015 \\
      LSST-Y10 & \texttt{lens} & \(m=10\) & -0.001 & +0.002 & -0.001 & +0.001 & +0.002 & +0.001 & +1.049 & +0.015 \\
      \midrule
      LSST-Y10 & \texttt{source} & \(m=1\) & -0.003 & +0.002 & -0.007 & +0.004 & +0.003 & +0.004 & +1.065 & +0.023 \\
      LSST-Y10 & \texttt{source} & \(m=2\) & -0.001 & +0.002 & +0.000 & +0.001 & +0.002 & +0.004 & +0.994 & +0.031 \\
      LSST-Y10 & \texttt{source} & \(m=3\) & +0.001 & +0.002 & -0.001 & +0.002 & -0.002 & +0.004 & +1.015 & +0.020 \\
      LSST-Y10 & \texttt{source} & \(m=4\) & -0.012 & +0.003 & -0.023 & +0.005 & +0.026 & +0.003 & +1.385 & +0.036 \\
      LSST-Y10 & \texttt{source} & \(m=5\) & -0.041 & +0.004 & -0.040 & +0.002 & +0.112 & +0.009 & +1.390 & +0.024 \\
      \bottomrule
    \end{tabular}
  \end{threeparttable}
\end{table*}

The tabulated values recapitulate in compact form the qualitative patterns documented in Sections~\ref{sec:4} and~\ref{sec:5}, whilst providing the precise numerical priors adopted in the released bias-corrected products of Section~\ref{sec:5.3}. For both LSST Y1 and Y10, lens-bin calibration is tightly controlled: the within-experiment scaled biases \(\bar{\delta}_{\mu_m}\) and \(\bar{\delta}_{\eta_m}\) remain at or below the \(10^{-3}\) level across every lens bin except the highest-redshift one of each scenario, where mild residual broadening emerges. Source-bin calibration is comparatively demanding, with residual biases concentrated in the highest-redshift bins: for Y1, the largest residual is in the fifth source bin at \(\bar{\delta}_{\mu_m}\simeq-0.016\); for Y10, the last two source bins (\(m=4,5\)) exceed the \(10^{-2}\) level, reaching \(\bar{\delta}_{\mu_m}\simeq-0.041\) at \(m=5\). These Y1 \(m=5\) and Y10 \(m=4,5\) source bins lie clearly outside the \(2\times 10^{-3}\) statistical-only SRD target quoted in Section~\ref{sec:4.3}, identifying high-redshift tail recovery as the dominant contribution to the residual calibration error budget. The corresponding experiment-to-experiment scatters \(\sigma_{\mu_m}\) and \(\sigma_{\eta_m}\) are also largest in these same bins, consistent with the elevated sensitivity to spectroscopic-selection and augmentation realisations identified in Section~\ref{sec:4.4}.

The experiment-marginalised block of the table lists the deterministic corrections \(\langle\Delta_m\rangle\) and \(\langle\Gamma_m\rangle\) that recentre each Dirichlet-sampled ensemble onto the supported \texttt{Truth} benchmark through the moment-matching transformations of Equations~(\ref{eq:36})--(\ref{eq:37}), together with their Dirichlet-ensemble one-sigma scatters \(\varsigma_{\Delta_m}\) and \(\varsigma_{\Gamma_m}\) defined in Section~\ref{sec:5.3}. Unlike the scaled biases \(\bar{\delta}_{\mu_m}\), which carry an explicit \(1/(1+\mu_m^{\texttt{Truth}})\) normalisation in their definition, the mean-shift parameter \(\langle\Delta_m\rangle\) is quoted directly in physical redshift units without any redshift-dependent scaling; the two therefore differ by a factor of \((1+\langle\mu_m^{\texttt{Truth}}\rangle)\), with \(\langle\Delta_m\rangle \approx -\bar{\delta}_{\mu_m}\,(1+\langle\mu_m^{\texttt{Truth}}\rangle)\) at leading order. For example, the Y10 \(m=5\) source bin with \(\bar{\delta}_{\mu_m}=-0.041\) yields \(\langle\Delta_m\rangle=+0.092\), consistent with an effective bin-centre redshift \(\langle\mu_m^{\texttt{Truth}}\rangle\) of order unity. The width-rescaling factors \(\langle\Gamma_m\rangle\) are multiplicative rather than additive, so values close to unity indicate no correction: they sit within a few per cent of unity across every lens bin, and rise to \(\langle\Gamma_m\rangle\simeq 1.4\text{--}1.5\) for the Y10 source bins \(m=4,5\), meaning that a \(40\text{--}50\%\) upward rescaling of the ensemble-mean width is needed to match \texttt{Truth} in those regimes. The residual scatters \(\varsigma_{\Delta_m}\) and \(\varsigma_{\Gamma_m}\) are small in absolute terms and remain at least an order of magnitude below the central offsets in the high-redshift source bins where the corrections are largest, confirming that the deterministic shifts are resolved robustly against experiment-marginalised fluctuations in the regimes where they matter most. Forecasting studies adopting the released priors should therefore treat \(\langle\Delta_m\rangle\) and \(\langle\Gamma_m\rangle\) as the per-bin deterministic offsets, \(\varsigma_{\Delta_m}\) and \(\varsigma_{\Gamma_m}\) as the corresponding one-sigma prior widths, and retain the cross-bin covariance structure shown in Figures~\ref{fig:8}--\ref{fig:9} whenever bin-to-bin correlations matter.
